\chapter{C-Reactive Protein (CRP)}

\section*{What Is This Test?}
The C-reactive protein (CRP) test measures the concentration of CRP in the blood. CRP is an acute-phase reactant synthesized primarily by hepatocytes in response to inflammatory cytokines (mainly IL-6, with contributions from IL-1 and TNF-alpha). It is a member of the pentraxin protein family (the pentraxin family also includes serum amyloid P). CRP binds to phosphocholine on the surface of damaged cells, bacteria, and fungi, and activates the complement system through C1q. It is a sensitive but nonspecific marker of acute inflammation, tissue injury, and infection. CRP levels rise within 4--6 hours of the inflammatory stimulus, double every 8 hours, peak at 24--48 hours, and fall rapidly with resolution (half-life of approximately 19 hours). Because of its rapid rise and fall, CRP is useful for detecting acute inflammation, monitoring response to therapy, and screening for occult infection. A highly sensitive CRP assay (hs-CRP) measures very low CRP levels (0.1--10 mg/L) and is used for cardiovascular risk stratification.

\section*{Alternative Names}
\begin{itemize}
  \item CRP
  \item High-Sensitivity C-Reactive Protein
  \item hs-CRP
  \item C-Reactive Protein, Quant
  \item Cardiac C-Reactive Protein
  \item Ultra-Sensitive CRP
  \item CRP (Immunoturbidimetric)
  \item Inflammation Marker Panel
\end{itemize}
\section*{Principle}
CRP is measured by immunoturbidimetric, immunonephelometric, or enzyme-linked immunosorbent assay (ELISA) methods. All use specific antibodies against human CRP. In immunoturbidimetry, CRP-antibody complexes scatter light, and the degree of light scatter is proportional to CRP concentration. The hs-CRP assay is a high-sensitivity version of the same immunoassay, with a lower limit of detection of approximately 0.1--0.3 mg/L (vs. 3--5 mg/L for the standard assay). hs-CRP is used for cardiovascular risk stratification: CRP <1.0 mg/L = low cardiovascular risk, 1.0--3.0 mg/L = intermediate risk, >3.0 mg/L = high risk. Serum CRP is measured in mg/L. In acute inflammation or infection, CRP may exceed 100 mg/L within 24--48 hours. In severe bacterial infections (sepsis), levels may reach 300--500 mg/L. CRP has a rapid response time, making it superior to the ESR for monitoring acute changes in inflammatory status.

\section*{History}
CRP was discovered in 1930 by William Tillett and Thomas Francis at Rockefeller University. It was so named because it reacted with the C-polysaccharide (C-carbohydrate) of the pneumococcus (Streptococcus pneumoniae) \cite{tillett1930serological}. It was later shown to be a protein that appears in the serum during acute illness, establishing it as an acute-phase protein. In 1950, it was found that CRP is produced in the liver. In the 1980s--1990s, it became clear that CRP is a highly sensitive marker of inflammation and that chronic low-grade inflammation plays a role in atherosclerosis. The JUPITER trial (Justification for the Use of Statins in Prevention: an Intervention Trial Evaluating Rosuvastatin, 2008) demonstrated that rosuvastatin reduced cardiovascular events in patients with elevated hs-CRP (>2.0 mg/L) but normal LDL cholesterol \cite{ridker2008rosuvastatin}. This established hs-CRP as a clinically useful marker of cardiovascular risk \cite{pepys2003creactive}.

\section*{Why This Test Is Ordered}
\begin{itemize}
  \item Assessment of acute inflammation or infection: fever of unknown origin, suspected sepsis, acute pancreatitis, inflammatory bowel disease flare, or postoperative infection. CRP >100 mg/L strongly suggests a bacterial infection
  \item Monitoring response to therapy in inflammatory diseases: rheumatoid arthritis, inflammatory bowel disease, vasculitis, SLE (CRP is less useful in SLE, where it often remains low even during flares), polymyalgia rheumatica, and giant cell arteritis (CRP is usually elevated in these conditions and falls with steroids)
  \item Cardiovascular risk stratification (hs-CRP): in asymptomatic patients with intermediate Framingham risk, a persistently elevated hs-CRP (>2.0 mg/L) identifies a higher-risk group who may benefit from statin therapy
  \item Evaluation of tissue injury: myocardial infarction (CRP peaks at 24--48 hours after onset), acute pancreatitis, burns, trauma, surgery (CRP rises postoperatively, peaking at 48--72 hours; a failure to fall may indicate a complication)
  \item Screening for occult infection in immunocompromised or neutropenic patients, and monitoring of antibiotic response (CRP doubling every 8 hours is characteristic of bacterial sepsis)
\end{itemize}

\section*{Primary vs Secondary Test}
CRP is a primary test for detecting acute inflammation and monitoring response to therapy. hs-CRP is used as a secondary/supplementary test for cardiovascular risk assessment, added to the standard lipid panel.

\section*{How This Test Is Performed}
A healthcare professional draws blood from a vein into a serum separator tube (gold-top) or lithium heparin tube (green-top). The sample is centrifuged and analyzed by immunoturbidimetric or nephelometric methods on automated chemistry analyzers. Results are typically available within 1--2 hours. Serum or plasma (heparin) may be used.

\section*{What Preparation Is Needed}
No fasting is required for a standard CRP. For hs-CRP, no fasting is required, but the sample should be drawn in a stable state: patients should not have an acute infection, recent surgery, or recent trauma within the preceding 2--3 weeks, as these cause transient elevations that confound cardiovascular risk assessment. Two hs-CRP measurements taken 2 weeks apart are recommended for cardiovascular risk stratification (the average is used). NSAIDs and statins may lower CRP; their use should be documented.

\section*{Contraindications and Precautions}
\begin{itemize}
  \item CRP is nonspecific: it rises in any inflammatory state (infection, autoimmune disease, malignancy, trauma, surgery, burns). It cannot by itself identify the cause of inflammation
  \item hs-CRP and standard CRP are different assays. Standard CRP is reported in mg/L (or mg/dL), with a normal range up to 3--5 mg/L. hs-CRP is reported in mg/L, with cardiovascular risk categories of <1.0 (low), 1.0--3.0 (intermediate), and >3.0 mg/L (high). They are not interchangeable for cardiovascular risk assessment
  \item CRP is not useful for diagnosing or monitoring SLE disease activity: CRP often remains normal or only mildly elevated during SLE flares, except in serositis (pleuritis/pericarditis), where CRP may rise
  \item CRP is elevated in pregnancy (physiologic, up to 20--30 mg/L), obesity (baseline CRP is 1.5--2x higher), and chronic kidney disease
  \item Steroids rapidly suppress CRP (within 24 hours) and may mask acute inflammation. Document steroid use when interpreting results
  \item In neonatal sepsis, CRP may be normal in the first 24--48 hours (slow hepatic response); a single negative CRP does not rule out sepsis, and serial measurements are recommended
\end{itemize}
\section*{Risks}
Standard venipuncture risks: mild pain, bruising, or hematoma at the puncture site.

\section*{What Happens After the Test}
Results are available within 1--2 hours. If CRP is elevated, the provider investigates the underlying cause with clinical assessment, additional laboratory tests, and imaging. If CRP was ordered for monitoring, the trend is more informative than a single value: a rising CRP over 48 hours suggests worsening inflammation or infection; a falling CRP suggests resolution or response to therapy. hs-CRP results for cardiovascular risk are interpreted with the lipid profile and the patient's overall risk factors.

\section*{Understanding Results}

\subsection*{Below Normal}
\begin{itemize}
  \item \textbf{CRP <0.3 mg/L (or undetectable):} Normal, indicating no significant acute inflammation. Most healthy individuals have CRP below 1.0 mg/L. In an adult without acute illness, a CRP of 0.1 mg/L is typical
  \item \textbf{CRP <0.5 mg/L on hs-CRP assay:} Low cardiovascular risk (<1.0 mg/L category). No additional hs-CRP-based risk modification is recommended
  \item \textbf{Low CRP with clinical signs of infection:} Possible early-phase (within 24 hours of onset) inflammatory response, steroid use, or an alternative explanation for the clinical findings. Repeat CRP in 12--24 hours if the clinical suspicion remains
  \item \textbf{CRP falling from a prior peak:} Indicates resolving inflammation or response to antibiotics/steroids. Serial monitoring is recommended
\end{itemize}

\subsection*{Above Normal}
\begin{itemize}
  \item \textbf{CRP 1.0--3.0 mg/L (hs-CRP):} Moderate cardiovascular risk. The patient is a candidate for enhanced risk factor modification (lifestyle, lipid control), particularly if other risk factors are present
  \item \textbf{CRP >3.0 mg/L (hs-CRP):} High cardiovascular risk. Per JUPITER trial criteria, statin therapy reduced cardiovascular events in this group. Exclude active infection or inflammation before attributing the elevation to vascular risk
  \item \textbf{CRP 10--100 mg/L:} Moderate inflammation. Seen in localized infection, autoimmune disease flares, acute myocardial infarction, pancreatitis, trauma, or major surgery. Interpretation requires clinical correlation
  \item \textbf{CRP >100 mg/L:} Marked elevation. Strongly suggests severe bacterial infection (sepsis, pneumonia, meningitis, peritonitis), acute pancreatitis, or major tissue necrosis. Viral infections typically produce lower CRP (<50 mg/L, often <20 mg/L)
  \item \textbf{CRP >300 mg/L:} Extreme elevation, seen in overwhelming sepsis, extensive burns, and acute severe pancreatitis. Critical illness is highly likely; prompt investigation and therapy are warranted
\end{itemize}

\section*{Normal Results}
\begin{itemize}
  \item \textbf{Adults (standard CRP assay):} <3.0--5.0 mg/L (values below 3.0 mg/L are normal; reference intervals vary by laboratory)
  \item \textbf{Adults (hs-CRP assay):} <1.0 mg/L is low cardiovascular risk; 1.0--3.0 mg/L is intermediate risk; >3.0 mg/L is high risk
  \item \textbf{Children:} <3.0 mg/L (age-dependent; slightly higher in the first year of life)
  \item \textbf{Pregnancy:} 20--30 mg/L (physiologic elevation); reference values are not established
\end{itemize}

\section*{Follow-up Tests}
\begin{itemize}
  \item \textbf{Erythrocyte sedimentation rate (ESR):} Both measure inflammation, but they are not interchangeable: CRP rises and falls rapidly (hours to days) and is a better marker of acute inflammation; ESR rises and falls slowly (days to weeks) and reflects chronic inflammation. They are frequently ordered together to characterize the inflammatory response
  \item \textbf{Complete blood count with differential:} Leukocytosis suggests bacterial infection; leukopenia with neutropenia suggests severe sepsis or hematologic disease. Combined with CRP, these help distinguish infection from other causes of inflammation
  \item \textbf{Blood cultures and procalcitonin:} When sepsis is suspected (CRP >100 mg/L with fever and hypotension). Procalcitonin rises more specifically in bacterial infection than CRP and may be used to guide antibiotic duration
  \item \textbf{Lipid panel and hs-CRP:} For cardiovascular risk stratification. hs-CRP is interpreted with LDL, HDL, and total cholesterol
  \item \textbf{Autoimmune panel:} ANA, anti-dsDNA, RF, anti-CCP, ANCA when autoimmune disease is suspected based on the clinical picture and persistent CRP elevation
  \item \textbf{Repeat CRP in 48--72 hours:} To monitor response to antibiotics or anti-inflammatory therapy. CRP falling by more than 50\% within 48 hours generally indicates a good response to therapy
\end{itemize}

\section*{References}
\bibliographystyle{unsrtnat}
\bibliography{references}

\clearpage
\section*{Regulatory Status and Authorization Date}
This chapter describes a test category, not a specific commercial product. FDA, CDSCO, EU IVDR, and MHRA approval or registration dates therefore cannot be assigned without the assay manufacturer, product name, instrument, and intended use. For laboratory-developed tests, an individual agency approval date may not exist. See the Preface for the interpretation of regulatory status.

\clearpage
